% lint-skip: 未完成ステータス（figure/table caption 未対応，人手チェック時に解消予定）
\documentclass[../../main.tex]{subfiles}
\begin{document}

\begin{flushright}
  \footnotesize\textit{【自動文字起こし・要確認】}
\end{flushright}

{\bf [解]}
$a_n > 0$
\[
  a_n^3 + 3 a_n^2 - \left(9 + \dfrac{1}{n}\right) a_n + 5 < 0 \quad \cdots \text{①}
\]
①を変形して
\[
  (a_n - 1)^2 < \dfrac{1}{n} \dfrac{a_n}{a_n + 5} \quad (\because a_n > 0) \quad \cdots \text{②}
\]
又，$f_n(x) = x^3 + 3x^2 - \left(9 + \dfrac{1}{n}\right) x + 5$ とおくと
\[
  f_n'(x) = 3x^2 + 6x - \left(9 + \dfrac{1}{n}\right)
\]
より，$f_n'(x) = 0$ の2根 $\alpha, \beta \ (\alpha < \beta)$ として下表を与える.

\begin{center}
  \begin{table}[htb]
    \centering
    \begin{tabular}{c|ccccc}
      $x$ & $\dots$ & $\alpha$ & $\dots$ & $\beta$ & $\dots$ \\ \hline
      $f_n'$ & $+$ & $0$ & $-$ & $0$ & $+$ \\ \hline
      $f_n$ & $\nearrow$ & & $\searrow$ & & $\nearrow$
    \end{tabular}
  \end{table}
\end{center}

\begin{figure}[htb]
  \centering
  \begin{tikzpicture}[x=1.15cm, y=0.16cm, >=stealth,
    every node/.style={font=\small}]
    % n=1 illustrates the sign pattern of the family f_n.
    \pgfmathsetmacro{\leftcrit}{-1-sqrt(4+1/3)}
    \pgfmathsetmacro{\rightcrit}{-1+sqrt(4+1/3)}
    \draw[->, thick, gray!80!black] (-3.9,0) -- (2.25,0) node[right] {$x$};
    \draw[->, thick, gray!80!black] (0,-2.5) -- (0,39) node[above] {$y$};
    \draw[line width=1.4pt, blue!85!black, domain=-3.7:1.9, samples=130]
    plot (\x,{\x*\x*\x+3*\x*\x-10*\x+5});
    \draw[dashed, gray!65] (\rightcrit,0)
    -- (\rightcrit,{\rightcrit*\rightcrit*\rightcrit+3*\rightcrit*\rightcrit-10*\rightcrit+5});
    \fill[red!85!black]
    (\rightcrit,{\rightcrit*\rightcrit*\rightcrit+3*\rightcrit*\rightcrit-10*\rightcrit+5}) circle (1.7pt);
    \node[below=4pt, font=\footnotesize] at (\leftcrit,0) {$\alpha$};
    \node[below=4pt, font=\footnotesize] at (\rightcrit,0) {$\beta$};
  \end{tikzpicture}
  \caption{$f_n$ の概形（図示には $n=1$ を使用）}
\end{figure}

従ってグラフは右上図のようになる. から，$\beta < x$ では $x$ は単調増加, $0 < x < \beta$ では $x$ は単調減少である. よって明らかに $0 < a_n < \dfrac{5}{3}$.

$y = \dfrac{x}{x+5}$ のグラフを考えてこの時

\begin{figure}[htb]
  \centering
  \begin{tikzpicture}[x=0.9cm, y=1.0cm, >=stealth,
    every node/.style={font=\small}]
    \draw[dashed, gray!65] (-5,-5) -- (-5,1.15);
    \draw[dashed, gray!65] (0,0.25) -- (4.4,0.25);
    \draw[dashed, gray!65] ({5/3},0) -- ({5/3},0.25);
    \draw[->, thick, gray!80!black] (-5.4,0) -- (4.7,0) node[right] {$x$};
    \draw[->, thick, gray!80!black] (0,-5) -- (0,1.2) node[above] {$y$};
    \draw[line width=1.4pt, blue!85!black, domain=-4.15:4.4, samples=130]
    plot (\x,{\x/(\x+5)});
    \fill[red!85!black] ({5/3},0.25) circle (1.6pt);
    \node[below=3pt, font=\footnotesize] at (-5,0) {$-5$};
    \node[below=3pt, font=\footnotesize] at ({5/3},0) {$\frac53$};
    \node[left=3pt, font=\footnotesize] at (0,0.25) {$\frac14$};
  \end{tikzpicture}
  \caption{$y=x/(x+5)$ と $x=5/3$ での値 $1/4$}
\end{figure}

\[
  0 < \dfrac{1}{n} \dfrac{a_n}{a_n + 5} < \dfrac{1}{n} \cdot \dfrac{1}{4}
\]
だから②より
\[
  (a_n - 1)^2 < \dfrac{1}{4n}
\]
したがってはさみうちの原理から
\[
  a_n \to 1 \quad (n \to \infty)
\]

\newpage

\end{document}
